\documentclass[10pt,a4paper]{amsart}
\usepackage[margin=19mm]{geometry}
\usepackage{amsmath,amssymb,mathtools}
\usepackage[scr=rsfs,cal=euler]{mathalfa}
\usepackage{xfrac}
\usepackage[unicode,hidelinks,bookmarks=false]{hyperref}
\newcommand{\ie}{i.e.\ }
\newcommand{\cf}{cf.\ }
\newcommand{\resp}{respectively\ }
\newcommand{\Subsection}{\subsection}
\providecommand{\nobreakdash}{\nobreak\hbox{-}}
\setlength{\emergencystretch}{2em}
\multlinegap0pt
\usepackage{amssymb}
\newtheorem{prop}{Proposition}
\title{On the Oja-Flow-Based Low-Rank Approximation of Kalman-Bucy Filters for Linear Time-Varying Systems}
\author{Kentaro Ohki}
\date{Source: arXiv:2607.29034v1 (CC BY 4.0)}
\begin{document}
\maketitle

\noindent\emph{Source and licence.} On the Oja-Flow-Based Low-Rank Approximation of Kalman-Bucy Filters for Linear Time-Varying Systems. Version arXiv:2607.29034v1 (CC BY 4.0), distributed by arXiv under the Creative Commons Attribution 4.0 International licence, http://creativecommons.org/licenses/by/4.0/, as recorded in the arXiv licence metadata for this exact version and shown on the abstract page https://arxiv.org/abs/2607.29034v1. Full licence notice: \texttt{licenses/arxiv-2607.29034-CC-BY-4.0.txt}.\par\smallskip

\noindent\textit{Editorial scope.} Counted unit: the article's prop prop:rotating\_A (source0.tex lines 152--181). The article's complete original proof of it is delivered with it (source0.tex lines 183--280, one \texttt{proof} environment of 98 lines). No mathematics is written, repaired or re-derived: the statement and the proof are the article's own lines, in the article's own notation, and no retained line is substituted or re-flowed.  Retained uncounted as the source-local context the proof needs: lines 100--103.  Nothing was omitted from any retained span, so the package carries no omission record.  No retained line contains a citation, so the wrapper carries no bibliography and no bibliography entry.  No retained line contains a figure environment, an image-inclusion command or drawing code, so no image is delivered and the package declares no asset dependency.  Editorial formatting only, in addition: an AMS \texttt{amsart} preamble that restates the article's own declarations and macros used by the retained lines, namely the environments \texttt{prop} on separate counters, together with the standard packages those lines need.  The retained environments print in this document as Proposition 1 in order of appearance, whereas the article prints the same items with the numbering of its own full document.  The complete native source of the article as distributed in the arXiv e-print for this exact version is bundled unchanged as \texttt{sources/206457/source0.tex}.
    \begin{align}
    \varepsilon \frac{d}{dt} U(t) = & ( I_{n} - U(t) U(t) ^{\top} ) A U(t) .
    \label{eq:Oja_flow}
    \end{align}

    \begin{prop}\label{prop:rotating_A}
    Let $(r,n) = (1,2)$ and consider the Oja flow \eqref{eq:Oja_flow} with the following time-varying matrix. 
    \begin{align*}
    A (t) = &
    R(t)
    \Lambda 
    R(t)^{\top},
    \\
    R(t) :=& \begin{bmatrix}
    \cos (\omega t) & \sin (\omega t) \\ -\sin (\omega t) & \cos (\omega t)
    \end{bmatrix},
     \ 
     \Lambda 
   := \begin{bmatrix}
    \lambda _{1} & 0 \\ 0 & \lambda _{2}, 
    \end{bmatrix}, 
    \end{align*}
    where $\lambda _{1} > \lambda _{2}$ and $\omega >0$. 
    Then, for any $\bar{\delta} \in (0,1)$, there exists $\varepsilon \in (0,1]$ such that $U(t)$ remains in the neighborhood of the dominant eigenvector $\psi _{1}(A(t)) =  \begin{bmatrix} \cos (\omega t) & -\sin (\omega t) \end{bmatrix}^{\top}$
    \begin{align*}
    & {\mathcal{N}}_{\bar{\delta}}( \psi _{1}(A(t))) 
    \\
    := & \big{\{} X \in {\mathrm{St}}(1,2) \ \big{|} \ 
    X = \psi _{1}(A(t))  + Y, 
    \\ 
    & \hspace{3cm} Y \in {\mathbb{R}}^{2}, \| Y \| _{\rm ind}< \bar{\delta}
    \big{\}}
    \end{align*}
    for $t\geq t_{0} \geq 0$ if $U(t_{0}) \in {\mathcal{N}}_{\delta }( \psi _{1}(A(t_{0})))$ where $\delta \in (0, \bar{\delta}/2)$. 
    \end{prop}

    \begin{proof}
    To move on the rotating frame, we consider the differential equation of $K(t) := R(t) ^{\top} U(t)$, 
    \begin{align*}
    &\varepsilon \frac{d}{dt} K(t) 
    \\
    =&
    \varepsilon \omega 
    \begin{bmatrix}
    0 & - \cos (2\omega t) \\ \cos (2\omega t) & 0
    \end{bmatrix} K(t)
    \\
    &
    +
    (I_{2} - K(t) K(t)^{\top}) \Lambda K(t)
    .
    \end{align*}
    Hence, the first term can be regarded as a perturbation or small modeling error. 
    Notice that if $\omega =0 $, then ${\mathcal{U}}_{1} := {\mathcal{U}}_{1}({\mathrm{diag}}(\lambda_{1}, \lambda _{2})) = \{ \begin{bmatrix} \pm 1 & 0 \end{bmatrix} ^{\top} \}$ is the stable equilibrium set. 
    However, ${\mathcal{U}}_{1}$ is no longer the equilibrium set when $R(t)$ is time-varying. 
    Since the second term of the right-hand side of the above equation forces the solution to ${\mathcal{U}}_{1}$, consider the linearization around an element of ${\mathcal{U}}_{1}$. 
    Let $K(t) = \bar{K} + K_{\delta}(t)$, where $\bar{K} \in {\mathcal{U}}_{1}$ and $K_{\delta}(t_{0}) \in {\mathcal{B}}_{\delta} := \{ X \in {\mathbb{R}}^{2\times 1} \ | \ \| X \| _{\rm ind} < \delta \ \}$ for a small positive $\delta \in (0, \bar{\delta} /2)$ at some $t_{0}\geq 0$. 
    Then, for $t\geq t_{0}$, 
    \begin{align*}
    & \varepsilon \frac{d}{dt} K_{\delta}(t) 
    \\
    =&
    -\varepsilon \omega \begin{bmatrix}
    0 & - \cos (2\omega t) \\ \cos (2\omega t) & 0
    \end{bmatrix} ( \bar{K} + K_{\delta}(t) )
    \\
    & +
    ( I_{2} - \bar{K} \bar{K}^{\top} ) \Lambda K_{\delta}(t)
    \\
    &
    -
    (\bar{K} K_{\delta} (t)^{\top} + K_{\delta}(t) \bar{K}^{\top} ) \Lambda \bar{K}
%    \\ & 
    + O (\delta ^2) 
    \end{align*}
    and without loss of generality, take $\bar{K} = \begin{bmatrix} 1 & 0 \end{bmatrix} ^{\top}$.  
    Then, the element-wise equation is 
    \begin{align*}
    \varepsilon \frac{d}{dt} K_{\delta , 2}(t) =&
     \varepsilon \omega (1+K_{\delta , 1}(t))
     \cos (2\omega t)
     \\ & 
    -(\lambda _{1} - \lambda _{2})
     K_{\delta,2}(t)
%     \\ &
    + O (\delta ^2) 
    \end{align*}
    and its solution is
    \begin{align*}
    K_{\delta,2} (t) = & \exp \left( -\frac{\lambda _{1} - \lambda _{2}}{\varepsilon}(t-t_{0}) \right) K_{\delta,2}(t_{0})
    \\ &
    +
    \omega \int _{t_{0}}^{t} \exp \left( -\frac{\lambda _{1} - \lambda _{2}}{\varepsilon} (t- \tau ) \right) 
    \\ & \hspace{1cm}
    \times (1+ K_{\delta ,1}(\tau )) \cos (2\omega \tau ) d\tau 
    \\ & 
    + O (\delta ^2  ) 
    .
    \end{align*}
    Therefore, 
    \begin{align*}
    | K_{\delta,2} (t) | \leq &  \exp \left( -\frac{\lambda _{1} - \lambda _{2}}{\varepsilon}(t-t_{0}) \right) |K_{\delta,2}(t_{0}) |
    \\ & 
    +
    \omega (1 + c(t) ) \int _{t_{0}}^{t} \exp \left( -\frac{\lambda _{1} - \lambda _{2}}{\varepsilon}\tau \right) d\tau 
    \\ & + O (\delta ^2  ) 
    \\
    \leq &
    \exp \left( -\frac{\lambda _{1} - \lambda _{2}}{\varepsilon}(t-t_{0}) \right) \delta 
    \\ & 
    +
    \frac{(1 + c(t) )\varepsilon \omega }{\lambda _{1} - \lambda _{2}} 
    \Bigg{(} 
    \exp \left( -\frac{\lambda _{1} - \lambda _{2}}{\varepsilon}t_{0} \right)
    \\ & \hspace{2cm}
    -
    \exp \left( -\frac{\lambda _{1} - \lambda _{2}}{\varepsilon}t \right)
    \Bigg{)}
    \\ & 
    + O (\delta ^2  ) 
    \\
    \leq &
    \exp \Bigg{(} -\frac{\lambda _{1} 
    - \lambda _{2}}{\varepsilon}(t-t_{0}) \Bigg{)} \delta 
    \\ &
    +
    \frac{(1+c(t))\varepsilon \omega }{\lambda _{1} - \lambda _{2}} 
    + O (\delta ^2  ) ,
    \end{align*}
    where $c(t) := \max _{\tau \in [t_{0},t]} |K_{\delta ,1}(\tau)|$. 
    Notice that $c_{\infty} := \sup _{t\geq t_{0}} c(t) < \infty$ because $\Psi(A(t))=R(t)$ is bounded and $U(t) \in {\mathrm{St}}(1,2)$. 
    Therefore, if for any $\delta >0$, taking $\varepsilon  \in (0,1]$ such that $\varepsilon  < \min \{ \delta , \delta (\lambda _{1} - \lambda _{2})/\omega (1+c_{\infty}) \}$ ensures the solution $ K_{\delta}(t) \in {\mathcal{B}}_{\bar{\delta}}$ for $t\geq t_{0}$. 
    The dominant mode is then $R(t) \bar{K} = \pm \begin{bmatrix} \cos (\omega t) & -\sin (\omega t) \end{bmatrix} ^{\top}$ and the solution remains near the dominant mode. 
    \end{proof}

\end{document}
